\section{DIFERENSIASI DAN ANALISIS PASAR}

\subsection{Diferensiasi Solusi}

Savora memiliki beberapa keunggulan diferensiasi yang membedakannya dari solusi \textit{food waste} yang sudah ada:

\begin{enumerate}[leftmargin=*]
    \item \textbf{Pendekatan Terpadu} — Berbeda dengan platform yang hanya fokus pada satu aspek (misalnya hanya penjualan diskon atau hanya donasi), Savora mengintegrasikan seluruh rantai nilai \textit{food waste} mulai dari penilaian kelayakan makanan, penjualan surplus dengan harga dinamis, hingga pelacakan dampak lingkungan \cite{kusumawati2025mobile}\cite{mobilefoodwaste2024review}.
    
    \item \textbf{Food Eligibility Assessment} — Sistem penilaian kelayakan makanan yang objektif dan transparan, memberikan kepastian bagi konsumen mengenai kualitas dan keamanan makanan yang dibeli. Menggunakan \textit{weighted scoring algorithm} berbasis 4 parameter utama.
    
    \item \textbf{Dynamic Pricing Engine} — Mesin penetapan harga dinamis yang mengoptimalkan harga berdasarkan sisa waktu kelayakan, permintaan pasar, dan margin UMKM \cite{qlearning2026pricing}\cite{kayikci2024iot}, memaksimalkan peluang penjualan sebelum makanan kadaluarsa.
    
    \item \textbf{Transparansi \& Trust Building} — Sistem \textit{Food Trust Index} yang menampilkan rating dan ulasan konsumen terhadap UMKM, serta \textit{Carbon Impact Tracking} yang menunjukkan dampak positif setiap pembelian terhadap lingkungan.
\end{enumerate}

\subsection{Keunggulan Teknis}

\subsubsection{Food Eligibility Score}
Sistem ini menggunakan model \textit{scoring} berbasis aturan yang mempertimbangkan:

\begin{itemize}[leftmargin=*]
    \item Waktu produksi dan estimasi masa simpan
    \item Kondisi kemasan dan suhu penyimpanan
    \item Riwayat penanganan makanan
    \item Kategori makanan dan karakteristik spesifik
\end{itemize}

\subsubsection{Dynamic Pricing Engine}
Mesin penetapan harga dinamis yang mengoptimalkan:

\begin{itemize}[leftmargin=*]
    \item Harga berdasarkan sisa waktu kelayakan
    \item Elasticitas permintaan lokal
    \item Harga kompetitor di area yang sama
    \item Margin minimum yang menguntungkan UMKM
\end{itemize}

\subsection{Technology Stack yang Unik}

Kombinasi teknologi yang digunakan dalam Savora memberikan keunggulan kompetitif:

\begin{table}[H]
\centering
\caption{Technology Stack Savora}
\label{tab:tech_stack}
\begin{tabularx}{\textwidth}{|l|Y|Y|}
\hline
\textbf{Komponen} & \textbf{Teknologi} & \textbf{Alasan Pemilihan} \\
\hline
Frontend & React.js & Komponen reusable, performa tinggi \\
\hline
Backend & Node.js + Express & asynchronous, cocok untuk real-time \\
\hline
Database & PostgreSQL + Redis & ACID compliance + caching \\
\hline
ML & Python + scikit-learn & Ekosistem ML yang mature \\
\hline
Real-time & Socket.IO & Komunikasi dua arah \\
\hline
Hosting & AWS / GCP & Skalabilitas dan reliability \\
\hline
\end{tabularx}
\end{table}

\subsection{Moat (Penghalang Kompetisi)}

Savora memiliki beberapa \textit{moat} yang melindungi posisi kompetitif:

\begin{enumerate}[leftmargin=*]
    \item \textbf{Data Lokal} — Dataset penjualan UMKM kuliner Indonesia yang dikumpulkan dari mitra merupakan aset berharga yang sulit ditiru oleh kompetitor baru atau perusahaan besar.
    
    \item \textbf{Jaringan UMKM} — Hubungan langsung dengan UMKM kuliner lokal memberikan keunggulan dalam hal \textit{network effect} dan \textit{switching cost}.
    
    \item \textbf{Algoritma yang Dioptimasi} — Model \textit{machine learning} yang dilatih dengan data lokal memiliki akurasi yang lebih tinggi dibandingkan model generik.
    
    \item \textbf{Mekanisme Donasi} — Integrasi dengan jaringan donasi memberikan nilai tambah yang sulit ditiru oleh platform yang hanya fokus pada penjualan.
\end{enumerate}
